\documentclass{article}
\usepackage[UTF8]{ctex}
\usepackage{amsmath,amssymb,amsfonts}
\usepackage{enumitem}
\usepackage{geometry}
\geometry{a4paper, margin=1in}
% === 错题本排版优化（不参与题目标识） ===
\setlength{\parindent}{0pt}
\setlength{\parskip}{0.35em plus 0.10em minus 0.08em}
\setlength{\emergencystretch}{3em}
\setlist{itemsep=0.25em, topsep=0.40em, parsep=0pt}
\allowdisplaybreaks[2]
% === 错题本排版优化结束 ===
\begin{document}
\section*{西南大学2024年数学分析真题}
\section*{64．西南大学2024年数学分析试卷}

\subsection{题目1}
1．函数极限的柯西收敛准则。

\subsection{题目2}
2．函数 \(f(x)\) 在区间 \(I\) 上不一致连续的定义．

\subsection{题目3}
3．函数极限 \(\lim _{x \rightarrow+\infty} f(x)\) 的归结原则。

\subsection{题目4}
1．求不定积分 \(\int \sin (\ln x) \mathrm{d} x\) ．

\subsection{题目5}
2 ．求极限 \(\lim _{x \rightarrow+\infty}\left[\frac{(1+x)^{\frac{1}{x}}}{e}\right]^{\frac{1}{x}}\) ．

\subsection{题目6}
3．设 \(f\) 具有二阶连续的偏导数，且 \(z=f\left(x^{2} y^{2}, x y\right)\) ，求 \(z_{x y}\) 。

三．（15 分）设有界函数 \(f(x)\) 在 \([a, b]\) 上的不连续点的全体为 \(\left\{x_{n} \mid n=1,2, \cdots\right\}\) ，且 \(\lim _{n \rightarrow \infty} x_{n}=a\) ，证明：\(f(x)\) 在 \([a, b]\) 上可积．\\
散．

四．（15 分）设数列 \(\left\{a_{n}\right\}\) 严格递减，且 \(a_{n}>0\) ．证明： \(\lim _{n \rightarrow \infty} a_{n}=0\) 的充要条件是级数 \(\sum_{n=1}^{\infty}\left(1-\frac{a_{n+1}}{a_{n}}\right)\)发

五．（15 分）已知 \(f(x)\) 在 \((-\infty,+\infty)\) 上二阶可导，且 \(f^{\prime \prime}(x) \leq 0, \lim _{x \rightarrow 0} \frac{f(x)}{2 x}=2\) ，证明：\(f(x) \leq 4 x\) ．

六．（15 分）用区间套定理证明确界原理．

七．（15 分）设函数 \(f(x)\) 在 \([1,+\infty)\) 上连续，恒正且单调递减，记

\[
u_{n}=\sum_{k=1}^{n} f(k)-\int_{1}^{n} f(x) \mathrm{d} x
\]

证明数列 \(\left\{u_{n}\right\}\) 收敛。

八．（15 分）设 \(f(x)\) 在 \([a,+\infty)\) 上一致连续，\(\varphi(x)\) 在 \([a,+\infty)\) 上连续，且 \(\lim _{x \rightarrow+\infty}[f(x)-\varphi(x)]=0\) ，证明：\(\varphi(x)\) 在 \([a,+\infty)\) 上一致连续。

九．（15 分）证明半径为 \(R\) 的球的体积为 \(\frac{4 \pi R^{3}}{3}\) ．

\subsection{题目14}
原题附加题：（15 分）计算曲线积分

\[
\int_{L} y \mathrm{~d} x+z \mathrm{~d} y+x \mathrm{~d} z
\]

其中 \(L\) 是 \(x^{2}+y^{2}+z^{2}=1\) 和 \(x+y+z=1\) 的交线，从 \(x\) 轴正向看去取逆时针方向．
\end{document}
